\section{Introduction and Motivation}

\subsection{Clinical Context: Thermal Ablation and Temperature Monitoring}

Thermal ablation has emerged as a minimally invasive treatment option for solid tumors across diverse anatomical sites, including the liver, lung, kidney, prostate, bone, and breast \cite{vogl2017, shibata2000, ahmed2014}. 
Among ablation modalities, Microwave Ablation (MWA) has gained significant clinical adoption due to several advantages over traditional radiofrequency ablation (RFA).
Compared to RFA, MWA produces faster heating rates, achieves higher intratumoral temperatures, creates larger ablation zones with shorter procedure times, demonstrates superior efficacy in high-perfusion areas due to reduced heat-sink effects, and generates more predictable and demarcated ablation zones \cite{vogl2017, lubner2010}.
These characteristics make MWA particularly suitable for treating large tumors (up to 5 cm) and lesions located adjacent to major blood vessels.

Successful thermal ablation requires precise control and monitoring to ensure complete tumor destruction while preserving healthy surrounding tissue.
Clinical standards define technically successful ablation as complete coverage of the target tumor plus a circumferential ablative margin of at least 5 mm (ideally 10 mm) \cite{ahmed2014}. 
Tissue coagulation and necrosis occur when temperature is maintained in the range of 50--60°C for at least 5 minutes; higher temperatures (>100°C) lead to vaporization and carbonization \cite{vogl2017}. 
This narrow therapeutic window necessitates continuous real-time monitoring of temperature dynamics to prevent both undertreatment (incomplete tumor necrosis leading to recurrence) and overtreatment (collateral damage to adjacent healthy structures).

Despite the clinical importance of temperature monitoring, direct measurement remains challenging. 
Invasive temperature sensors such as thermocouples and fiber-optic probes provide high temporal resolution but sample only discrete spatial locations, add procedural complexity, and increase infection risk. 
Therefore, clinicians seek non-invasive temperature imaging modalities.

\subsection{Current Monitoring Modalities: Capabilities and Limitations}

Non-invasive temperature monitoring for thermal ablation can be accomplished through several complementary imaging modalities, each with distinct trade-offs \cite{geoghegan2022}:

\begin{itemize}
    \item \textbf{Magnetic Resonance Thermometry (MRT)} is considered the gold standard for non-invasive temperature mapping, offering high spatial and temporal resolution and accurate quantitative temperature estimates. 
    However, MRT is prohibitively expensive, immobile, and requires complex MRI-compatible ablation equipment. Moreover, MRT is contraindicated in patients with pacemakers and ferromagnetic implants, and its availability remains limited in most clinical settings \cite{vogl2017, geoghegan2022}.

    \item \textbf{Computed Tomography (CT)} can delineate the ablation zone post-procedure through visualization of necrotic tissue but provides no real-time feedback during the procedure.
    Additionally, CT delivers ionizing radiation to both patient and clinician \cite{geoghegan2022}.

    \item \textbf{Ultrasound (US) B-mode imaging} is widely available, portable, cost-effective, non-invasive, and permits real-time observation during the procedure. 
    However, conventional B-mode ultrasound offers poor intrinsic contrast between the thermal lesion and surrounding normal tissue, limiting the ability to precisely demarcate ablation zone boundaries \cite{zhang2020, geoghegan2022}. 
    Indirect visual markers of heating such as hyperechogenic microbubbles, speckle decorrelation, and echo shift can be observed but do not correlate reliably with absolute tissue temperature \cite{seip2002}.

    \item \textbf{Invasive thermometry} (thermocouples, fiber-optic sensors) provides accurate point measurements at discrete locations but suffers from spatial limitations and procedural complexity \cite{geoghegan2022}.
\end{itemize}

Given these trade-offs, ultrasound-based monitoring is an attractive candidate for routine clinical deployment due to its ubiquity, real-time capability, and absence of radiation. 
However, the indirect relationship between observable ultrasound markers and absolute tissue temperature represents a significant scientific challenge.

\subsection{Deep Learning for Ultrasound-Based Monitoring}

Recent advances in deep learning have demonstrated the potential to enhance ultrasound-based monitoring of thermal lesions.
Zhang et al. \cite{zhang2020} proposed a Convolutional Neural Network (CNN) architecture for detecting and segmenting thermally induced lesions in ultrasound radiofrequency (RF) data acquired from porcine livers during MWA.
Their modified CNN achieved an area under the receiver operating characteristic curve (AUC) of 0.8948, substantially outperforming conventional B-mode image analysis (AUC 0.6904). This seminal work demonstrates that deep learning can extract features from ultrasound data that are more predictive of thermal damage than manual visual inspection.

However, existing deep learning approaches for ultrasound monitoring of ablation have focused primarily on lesion detection and segmentation---classifying whether tissue is coagulated or not---rather than regressing continuous temperature values. 
Furthermore, most state-of-the-art models are computationally intensive and designed for post-hoc analysis rather than real-time clinical deployment. 
There remains a significant gap between proof-of-concept studies and clinically feasible real-time systems.

\subsection{Physics-Informed Machine Learning}

An emerging paradigm in medical machine learning is the integration of domain knowledge and physical laws into the learning process. 
Physics-Informed Neural Networks (PINNs) represent one such approach, where partial differential equations (PDEs) governing physical processes are incorporated as soft constraints during training \cite{raissi2019physics}. 

For thermal ablation, the relevant physics is captured by the bioheat equation (Pennes equation), which governs temperature evolution in tissue:
\begin{equation}
    \frac{\partial T}{\partial t} = \beta \nabla^2 T - \alpha(T - T_{arterial}) + Q_{ext}
\end{equation}
where $T$ is tissue temperature, $\beta$ is thermal conductivity, $\alpha$ is the perfusion coefficient accounting for heat dissipation due to blood flow, $T_{arterial}$ is arterial blood temperature, and $Q_{ext}$ represents external heating from the ablation device.

By constraining neural networks to respect this equation, PINNs can improve generalization, particularly in data-scarce regimes common in medical imaging. 
The physical constraints act as a form of inductive bias, reducing the hypothesis space and improving robustness to noise. 
However, implementing PINNs for video regression from ultrasound sequences is non-trivial and has not been extensively studied for ablation monitoring.

\subsection{Edge Computing and Real-Time Deployment Requirements}

Clinical deployment of machine learning models in interventional radiology imposes stringent constraints:

\begin{itemize}
    \item \textbf{Data Privacy}: Transmitting sensitive patient imaging data to cloud servers for processing violates privacy regulations (HIPAA, GDPR) and introduces security risks. 
    Inference must occur locally on-device.
    
    \item \textbf{Real-Time Requirements}: High-Intensity Focused Ultrasound (HIFU) and interventional procedures require feedback at rates exceeding 30 FPS to enable timely operator response to thermal runaway or insufficient heating.
    
    \item \textbf{Hardware Constraints}: To enable retrofit to existing ultrasound machines and ensure accessibility, models must run on low-cost, resource-constrained devices such as the Raspberry Pi 4 or NVIDIA Jetson Nano, which feature limited RAM and CPU/GPU capacity.
    
    \item \textbf{Interpretability}: Black-box predictions are clinically unacceptable. Clinicians require confidence intervals and uncertainty estimates alongside temperature predictions to assess reliability and inform treatment decisions.
\end{itemize}

These constraints are often in tension: achieving high accuracy typically demands large, computationally intensive models, while edge deployment demands extreme efficiency. 
Furthermore, uncertainty quantification---critical for clinical safety---adds additional computational overhead.

\subsection{Motivation and Research Gaps}

Despite progress in deep learning for medical imaging and physics-informed learning, several research gaps remain:

\begin{enumerate}
    \item Most existing deep learning models for ultrasound monitoring perform classification (coagulated vs. not coagulated) rather than continuous temperature regression, which is necessary for precise ablation control.
    
    \item There is limited evidence that physics-informed constraints can improve video regression models for temperature estimation in ablation monitoring.
    
    \item Real-time deployment of video regression models on edge devices while maintaining accuracy has not been comprehensively benchmarked.
    
    \item Uncertainty quantification for temperature estimation in clinical settings remains underexplored; most thermal ablation studies report point estimates without confidence intervals.
    
    \item The role of temporal information (e.g., optical flow) as an explicit input feature to capture transient thermal dynamics has not been systematically investigated for ultrasound-based ablation monitoring.
\end{enumerate}

\subsection{Contributions of This Work}

To address these gaps, this paper presents a comprehensive study of deep learning architectures for real-time temperature regression from ultrasonic video sequences, optimized for edge deployment and augmented with principled uncertainty quantification.

Our main contributions are:

\begin{itemize}
    \item \textbf{Lightweight Architecture Design}: We propose a custom CNN-LSTM architecture with only $\sim$60k parameters, optimized for real-time inference on CPU-class hardware while maintaining competitive accuracy. 
    We compare this against pretrained ResNet-based models and physics-informed variants.
    
    \item \textbf{Physics-Informed Regularization}: We incorporate domain knowledge via physics-informed loss functions based on the bioheat equation, enabling the model to generate physically plausible predictions and improve generalization in data-limited regimes. 
    We explore both simple cooling constraints and advanced formulations incorporating tissue perfusion, spatial diffusion, convection, and metabolic heat generation.
    
    \item \textbf{Explicit Temporal Dynamics via Optical Flow}: We integrate dense optical flow as an explicit input feature to capture motion-induced temperature changes (e.g., convection currents, tissue deformation), enriching the model's ability to infer thermal dynamics from visual cues.
    
    \item \textbf{Uncertainty Quantification}: We implement two complementary probabilistic approaches---Deep Ensembles and Bayesian Neural Networks (BNNs)---to provide clinically actionable confidence intervals alongside temperature predictions. We propose a novel Bayesian Physics-Informed Neural Network (B-PINN) that synergizes uncertainty quantification with physical constraints.
    
    \item \textbf{Edge Deployment Evaluation}: We rigorously benchmark inference latency on CPU-constrained environments, demonstrating that optimized models achieve real-time performance (>30 FPS) on Raspberry Pi 4 and Jetson Nano, while maintaining high predictive accuracy (MAE $\le 1°$C).
    
    \item \textbf{Leave-One-Sequence-Out Validation}: We employ rigorous cross-validation strategies that account for temporal correlations within experimental runs, ensuring unbiased generalization estimates.
\end{itemize}

The remainder of this paper is organized as follows. Section \ref{sec:related_work} reviews related work on non-invasive temperature monitoring, deep learning in medical imaging, and physics-informed machine learning.
 Section \ref{sec:methodology} details our proposed architectures, preprocessing pipeline, and uncertainty quantification techniques. Section \ref{sec:experiments} describes the experimental setup, datasets, and evaluation metrics. Section \ref{sec:results} presents comparative results across architectures, and Section \ref{sec:discussion} contextualizes findings with respect to clinical deployment and future work.

